\documentclass[11pt,a4paper]{article}
\usepackage[utf8]{inputenc}\usepackage[T1]{fontenc}
\usepackage[provide=*,english]{babel}
\usepackage{lmodern,microtype,amsmath,amssymb,booktabs,geometry,xcolor,fancyhdr}
\usepackage[hidelinks]{hyperref}
\geometry{margin=2.3cm,headheight=14pt}
\setlength{\parindent}{0pt}\setlength{\parskip}{0.4em}
\definecolor{gcblue}{RGB}{34,73,110}\definecolor{gcgray}{RGB}{90,96,102}
\pagestyle{fancy}\fancyhf{}
\fancyhead[L]{\small\color{gcgray}GC2D: RK4}
\fancyhead[R]{\small\color{gcgray}Model theory}\fancyfoot[C]{\thepage}
\setlength{\emergencystretch}{3em}
\begin{document}
\begin{center}
 {\LARGE\bfseries\color{gcblue}Classical Fourth-Order Runge--Kutta}\par
 \vspace{0.35em}{\large Four-stage explicit integration}
\end{center}

\section{Butcher tableau and stages}

The classical method has tableau
\[
\begin{array}{c|cccc}
0&0&0&0&0\\
\tfrac12&\tfrac12&0&0&0\\
\tfrac12&0&\tfrac12&0&0\\
1&0&0&1&0\\ \hline
&\tfrac16&\tfrac13&\tfrac13&\tfrac16
\end{array}
\]
and, for $\dot z=f(t,z)$,
\begin{align}
k_1&=f(t_n,z_n),\\
k_2&=f(t_n+h/2,z_n+hk_1/2),\\
k_3&=f(t_n+h/2,z_n+hk_2/2),\\
k_4&=f(t_n+h,z_n+hk_3),\\
z_{n+1}&=z_n+\frac h6(k_1+2k_2+2k_3+k_4).
\end{align}
The non-autonomous stage times are part of the model.

\section{Order and stability}

The tableau satisfies the Runge--Kutta order conditions through order four;
the local error is $\mathcal O(h^5)$ and the global error is $\mathcal O(h^4)$
for sufficiently smooth problems.  On $y'=\lambda y$, its stability polynomial
is
\[
R(\xi)=1+\xi+\frac{\xi^2}{2}+\frac{\xi^3}{6}+\frac{\xi^4}{24}.
\]
The finite stability region does not make RK4 an A-stable method.

\section{Geometric properties}

Classical RK4 is not symmetric and is not symplectic for a general Hamiltonian
system.  Its high local accuracy can nevertheless make short-time energy and
area errors small.  The distinction is structural: symplecticity is the exact
identity
\[
(D\Phi_h)^T\Omega D\Phi_h=\Omega,
\]
not merely a small defect at one step size.  The implementation exposes each
complete step map to observers so local and accumulated defects can be
measured.

\section{Energy-tracking configuration}

An optional \texttt{EnergyObserver} is supplied to \texttt{simulate} through
\texttt{observers}. Execution binds the Hamiltonian functions and the observer
uses the retained RK4 stages to integrate the passive equation
\[
\dot k=-\partial_tH(t,z),\qquad K(z,t,k)=H(t,z)+k.
\]
The same RK4 nodes and weights are used for $z$ and the normalized momentum.  The physical
trajectory is returned normally; $k$ and the maximum sampled drift of $K$ are
diagnostics.  This extension monitors generalized energy but does not turn RK4
into a symplectic or energy-preserving method.

\section{Implementation correspondence}

The method is implemented in \texttt{src/methods/classical/rk4.py}
and consumes the injected vector field $f(t,z)$.  Main-grid steps are independent
of the output schedule; shadow steps provide off-grid samples without changing
the accepted trajectory or emitting observations.  See the sibling runtime
document for the exact result and diagnostic contract.

\section{Physical accepted states and passive energy tracking}
For $N$ particles with $d$ physical components, every accepted snapshot contains
$Z\in\mathbb{R}^{dN}$ and one scalar time $t$. An optional energy observer
owns separate physical momentum $\kappa\in\mathbb{R}^{N}$.
Thus GC accepted vectors have $2N$
components and FC vectors have $4N$, whether energy is tracked or not.
A doubled method prepares temporary $(u,v)\in\mathbb{R}^{2dN}$ copies,
allows them to separate, and projects to one physical result before acceptance.
The representation stores only the physical size and never owns a clock,
momentum, physical coordinate declaration, dynamics or solver.

The normalized passive balance is
\[
\dot\kappa=-\partial_tH(t,Z),\qquad \kappa(t_0)=0,\qquad
\Delta K(t)=H(t,Z(t))+\kappa(t)-H(t_0,Z_0).
\]
GC split-map contributions are summed in temporary work and divided by two
to obtain physical momentum. Tracking leaves physical stages, nonlinear
stopping scales and adaptive acceptance unchanged. An autonomous Hamiltonian
supplies an explicit zero momentum derivative.

Hairer's constraint is solely $u-v=0$: the reduced multiplier has $2N$ entries
and ABBA's simultaneous spatial solve has $6N$ unknowns. Neither $t$ nor
$\kappa$ is projected. The former eight-component projection is retired.
Dynamics declare component order and position/velocity roles through immutable
structural data; initial conditions and results retain those data independently.
Execution injects $f(t,Z)$ and optional physical derivatives into a fresh method
run. Methods own physical stages and projection equations, without retaining the
dynamics object. Nonlinear tolerances, iteration limits, differentiation and
iteration callbacks belong to the configured algebraic solver. Generic doubled
copy operations live in \path{src/formulations/state.py}; classical, HBVM and
adaptive methods need no state formulation.

Methods optionally return retained stage quadrature or a dense interpolant.
The energy observer integrates momentum independently, using Hamiltonian
functions supplied by execution. Controllers keep only physical accepted states;
compiled JAX loops carry the observer accumulator separately.
Observers receive physical states. Energy histories in \texttt{States.diagnostics}
retain their existing labels; \texttt{extended\_time} has shape
$(N,n_{\rm samples})$ and is broadcast from scalar sample times.
Fixed shadow samples leave accepted states and counters unchanged. Adaptive
quadrature uses accepted dense output and physical-only error control.

The passive momentum uses the same four stage states and weights as the physical RK4 update.

\section{CPU and JAX execution of the same map}

The execution configuration selects either the default NumPy/SciPy CPU driver
or a JAX driver on one CPU or GPU device. Both consume the same RK4 stage and
passive-quadrature formulas. The potential uses the same fitted spline knots
and coefficients in both implementations. JAX requires explicit double
precision; different floating-point operation ordering can produce round-off
differences, so equivalence does not mean bitwise identity.

For independent particles, stage operations are batched over the particle
index. The four stages retain their dependencies, and successive time steps
remain sequential inside a compiled loop. The accepted state and requested
sample buffer remain on the chosen device until the completed result is
transferred to the standard immutable NumPy solution. Independent shortened
maps still provide off-grid samples without altering accepted main states.
The passive momentum and sampled Hamiltonian are evaluated on that device.

The compiled route currently covers the built-in guiding-centre and
full-cyclotron equations. Custom dynamics, method subclasses and Python
step callbacks or progress reporting require the CPU route. Device choice
does not change the order, stability region or geometric properties derived
above. First-call compilation must be distinguished from warmed integration
timings, which include transfers and output construction in the test study.

\end{document}
